% GENERATED FILE. Edit canonical lessons, then run npm run generate:books.
% canonical-source-hash: fnv1a64:e4f39c824824667d
% canonical-lessons: RU-C273-proiskhodit, RU-C273-prodolzhatsya, RU-C273-dlitsya, RU-C273-prekrashchatsya, RU-C273-nastupat

\chapter{Happening and going on}
\label{ch:ru-happening-and-going-on}
\hlchaptermodality{273}

\begin{chapteropening}
\textbf{By the end of this chapter:} \emph{I can say that something happens, goes on, lasts, stops and sets in.}

Complete the last lesson of chapter 273: I can say that something happens, goes on, lasts, stops and sets in.
\end{chapteropening}

\section[\texorpdfstring{\ru{происходить} (proiskhodít)}{происходить (proiskhodít)}]{\ru{происходить} (proiskhodít) — to happen, to take place}

\label{lesson:RU-C273-proiskhodit}



\begin{quote}
Before the new one: say the Russian for soon after, then the Russian for later.
\end{quote}

\subsection*{You'll want to know: \ru{происходить}}
\textbf{\ru{происходить}} (\emph{proiskhodít}) — \textquotedblleft{}to happen, to take place\textquotedblright{}.

\begin{quote}
\ru{Это} \ru{происходит} \ru{каждую} \ru{зиму}.
\end{quote}

\begin{quote}
\emph{Éto proiskhódit kázhduyu zímu.}
\end{quote}

\begin{quote}
This happens every winter.
\end{quote}

Imperfective; the perfective is \emph{\ru{произойти}} (\emph{proizoytí}), past \emph{\ru{произошёл}} (\emph{proizoshyól}): \emph{\ru{что} \ru{произошло}?} — what happened?

\begin{cousinweb}
\emph{\ru{происходить}} is built on \emph{\ru{про}-}, \emph{\ru{из}-} and \emph{\ru{ходить}} \textquotedblleft{}to go\textquotedblright{}, which is inherited from Proto-Slavic \emph{xoditi} (EtymDB): to come out, to come about.
\end{cousinweb}

\begin{grammarlens}[title={What you've built}]
One more word for telling, explaining and describing at B1.
\end{grammarlens}

\subsection*{Your turn}
\begin{itemize}
\raggedright
  \item \emph{Say it:} \emph{proiskhodít}
  \item \emph{Say it:} \emph{proiskhodít}, once more
  \item \emph{Say it:} a sentence of your own with \emph{proiskhodít}
  \item \emph{Recall:} say the Russian for soon after, then the Russian for later, then say \emph{proiskhodít} again
\end{itemize}

\begin{tcolorbox}[breakable,colback=teal!4,colframe=teal!35!black,title={Before you move on}]
What does \ru{происходить} mean? (\textquotedblleft{}To happen, to take place\textquotedblright{}.) Say it once more.
\end{tcolorbox}
\section[\texorpdfstring{\ru{продолжаться} (prodolzhátsya)}{продолжаться (prodolzhátsya)}]{\ru{продолжаться} (prodolzhátsya) — to go on, to continue}

\label{lesson:RU-C273-prodolzhatsya}



\begin{quote}
Before the new one: say the Russian for later, then the Russian for to happen.
\end{quote}

\subsection*{You'll want to know: \ru{продолжаться}}
\textbf{\ru{продолжаться}} (\emph{prodolzhátsya}) — \textquotedblleft{}to go on, to continue\textquotedblright{}.

\begin{quote}
\ru{Дождь} \ru{продолжался} \ru{три} \ru{дня}.
\end{quote}

\begin{quote}
\emph{Dozhd prodolzhálsya tri dnya.}
\end{quote}

\begin{quote}
The rain went on for three days.
\end{quote}

Imperfective; perfective \emph{\ru{продолжиться}} (\emph{prodólzhitsya}). Without \emph{-\ru{ся}}, \emph{\ru{продолжать}} takes an object: to continue something.

\begin{cousinweb}
\emph{\ru{продолжаться}} is built on \emph{\ru{про}-} and \emph{\ru{долгий}} \textquotedblleft{}long\textquotedblright{}, which is inherited from Proto-Slavic \emph{dĭlgŭ} \textquotedblleft{}long\textquotedblright{} (EtymDB): to be made longer.
\end{cousinweb}

\begin{grammarlens}[title={What you've built}]
One more word for telling, explaining and describing at B1.
\end{grammarlens}

\subsection*{Your turn}
\begin{itemize}
\raggedright
  \item \emph{Say it:} \emph{prodolzhátsya}
  \item \emph{Say it:} \emph{prodolzhátsya}, once more
  \item \emph{Say it:} a sentence of your own with \emph{prodolzhátsya}
  \item \emph{Recall:} say the Russian for later, then the Russian for to happen, then say \emph{prodolzhátsya} again
\end{itemize}

\begin{tcolorbox}[breakable,colback=teal!4,colframe=teal!35!black,title={Before you move on}]
What does \ru{продолжаться} mean? (\textquotedblleft{}To go on, to continue\textquotedblright{}.) Say it once more.
\end{tcolorbox}
\section[\texorpdfstring{\ru{длиться} (dlítsya)}{длиться (dlítsya)}]{\ru{длиться} (dlítsya) — to last}

\label{lesson:RU-C273-dlitsya}



\begin{quote}
Before the new one: say the Russian for to happen, then the Russian for to go on.
\end{quote}

\subsection*{You'll want to know: \ru{длиться}}
\textbf{\ru{длиться}} (\emph{dlítsya}) — \textquotedblleft{}to last\textquotedblright{}.

\begin{quote}
\ru{Фильм} \ru{длился} \ru{два} \ru{часа}.
\end{quote}

\begin{quote}
\emph{Film dlílsya dva chasá.}
\end{quote}

\begin{quote}
The film lasted two hours.
\end{quote}

Imperfective; perfective \emph{\ru{продлиться}} (\emph{prodlítsya}). The time is in the accusative: \emph{\ru{урок} \ru{длится} \ru{час}}.

\begin{cousinweb}
\emph{\ru{длиться}} is built on \emph{\ru{длина}} \textquotedblleft{}length\textquotedblright{}, from the same Slavic root as \emph{\ru{долгий}} \textquotedblleft{}long\textquotedblright{}, Proto-Slavic \emph{dĭlgŭ} (EtymDB).
\end{cousinweb}

\begin{grammarlens}[title={What you've built}]
One more word for telling, explaining and describing at B1.
\end{grammarlens}

\subsection*{Your turn}
\begin{itemize}
\raggedright
  \item \emph{Say it:} \emph{dlítsya}
  \item \emph{Say it:} \emph{dlítsya}, once more
  \item \emph{Say it:} a sentence of your own with \emph{dlítsya}
  \item \emph{Recall:} say the Russian for to happen, then the Russian for to go on, then say \emph{dlítsya} again
\end{itemize}

\begin{tcolorbox}[breakable,colback=teal!4,colframe=teal!35!black,title={Before you move on}]
What does \ru{длиться} mean? (\textquotedblleft{}To last\textquotedblright{}.) Say it once more.
\end{tcolorbox}
\section[\texorpdfstring{\ru{прекращаться} (prekrashchátsya)}{прекращаться (prekrashchátsya)}]{\ru{прекращаться} (prekrashchátsya) — to stop, to cease}

\label{lesson:RU-C273-prekrashchatsya}



\begin{quote}
Before the new one: say the Russian for to go on, then the Russian for to last.
\end{quote}

\subsection*{You'll want to know: \ru{прекращаться}}
\textbf{\ru{прекращаться}} (\emph{prekrashchátsya}) — \textquotedblleft{}to stop, to cease\textquotedblright{}.

\begin{quote}
\ru{Музыка} \ru{не} \ru{прекращалась} \ru{всю} \ru{ночь}.
\end{quote}

\begin{quote}
\emph{Múzyka ne prekrashchálas vsyu noch.}
\end{quote}

\begin{quote}
The music did not stop all night.
\end{quote}

Imperfective; perfective \emph{\ru{прекратиться}} (\emph{prekratítsya}). Without \emph{-\ru{ся}}, \emph{\ru{прекращать}} stops something else.

\begin{cousinweb}
\emph{\ru{прекращаться}} is built on \emph{\ru{пре}-} and \emph{\ru{кратить}} \textquotedblleft{}to shorten\textquotedblright{}, from \emph{\ru{краткий}} \textquotedblleft{}short\textquotedblright{}, usually explained as the Church Slavonic shape of native \emph{\ru{короткий}}, Proto-Slavic \emph{kortŭkŭ} (EtymDB): to cut short.
\end{cousinweb}

\begin{grammarlens}[title={What you've built}]
One more word for telling, explaining and describing at B1.
\end{grammarlens}

\subsection*{Your turn}
\begin{itemize}
\raggedright
  \item \emph{Say it:} \emph{prekrashchátsya}
  \item \emph{Say it:} \emph{prekrashchátsya}, once more
  \item \emph{Say it:} a sentence of your own with \emph{prekrashchátsya}
  \item \emph{Recall:} say the Russian for to go on, then the Russian for to last, then say \emph{prekrashchátsya} again
\end{itemize}

\begin{tcolorbox}[breakable,colback=teal!4,colframe=teal!35!black,title={Before you move on}]
What does \ru{прекращаться} mean? (\textquotedblleft{}To stop, to cease\textquotedblright{}.) Say it once more.
\end{tcolorbox}
\section[\texorpdfstring{\ru{наступать} (nastupát)}{наступать (nastupát)}]{\ru{наступать} (nastupát) — to come, to set in (of a time); to step on}

\label{lesson:RU-C273-nastupat}



\begin{quote}
Before the new one: say the Russian for to last, then the Russian for to stop.
\end{quote}

\subsection*{You'll want to know: \ru{наступать}}
\textbf{\ru{наступать}} (\emph{nastupát}) — \textquotedblleft{}to come, to set in (of a time); to step on\textquotedblright{}.

\begin{quote}
\ru{Наступает} \ru{вечер}.
\end{quote}

\begin{quote}
\emph{Nastupáyet vécher.}
\end{quote}

\begin{quote}
Evening is coming on.
\end{quote}

Imperfective; perfective \emph{\ru{наступить}} (\emph{nastupít}): \emph{\ru{наступила} \ru{зима}} — winter came. Also \emph{\ru{наступить} \ru{на} \ru{ногу}} — to step on a foot.

\begin{cousinweb}
\emph{\ru{наступать}} is built on \emph{\ru{на}-} and \emph{\ru{ступать}} \textquotedblleft{}to step\textquotedblright{}, the root of \emph{\ru{ступень}} \textquotedblleft{}a step\textquotedblright{}, inherited from Proto-Slavic \emph{stǫpenĭ} (EtymDB): to step onto.
\end{cousinweb}

\begin{grammarlens}[title={What you've built}]
One more word for telling, explaining and describing at B1.
\end{grammarlens}

\subsection*{Your turn}
\begin{itemize}
\raggedright
  \item \emph{Say it:} \emph{nastupát}
  \item \emph{Say it:} \emph{nastupát}, once more
  \item \emph{Say it:} a sentence of your own with \emph{nastupát}
  \item \emph{Recall:} say the Russian for to last, then the Russian for to stop, then say \emph{nastupát} again
\end{itemize}

\begin{tcolorbox}[breakable,colback=teal!4,colframe=teal!35!black,title={Before you move on}]
What does \ru{наступать} mean? (\textquotedblleft{}To come, to set in (of a time); to step on\textquotedblright{}.) Say it once more.
\end{tcolorbox}
